\scanpage{66}

Is \inserted{such} an automorphism $u_0$ determined by its restriction to $\widehat\Pi_{0,3}$? In other words, an automorphism of $\widehat{\mathfrak S}_{0,3}^{+}$ \deleted{which} ($\simeq\mathrm{Sl}(\unsure{2|1},\mathbb Z)^{\unsure{\prime}\wedge}$) which induces the identity on $\widehat\Pi_{0,3}$ and on $\mathfrak S_{3}$, is it trivial? It suffices to prove that its restriction to the inverse image \inserted{$\widehat\Pi_{\unsure{0,\infty|0,3}}\cdot\{1,\unsure{\sigma_\infty}\}$} of $\{1,\dot\sigma_\infty\}\subset\mathfrak S_{3}$ is \unsure{trivial|trivial} (by symmetry this will also be the case on replacing $\dot\sigma_\infty$ by $\dot\sigma_0,\dot\sigma_1$, now these three subgroups generate $\widehat{\mathfrak S}_{0,3}^{+}$. We will have $f(\sigma_\infty)=\sigma_\infty g$, with $g\in\widehat\Pi_{0,3}$ [and we must have $\sigma_\infty g\sigma_\infty g=1$ i.e.
\[
g\cdot\sigma_\infty(g)=1
\]
]
moreover \deleted{for $x\in\Pi_{0,3}$ \unsure{we have}} \inserted{we will have to have that}
\[
\operatorname{int}(\sigma_\infty g)\vert\widehat\Pi_{0,3}
=\deleted{\unsure{f(x\sigma_\infty)}}\operatorname{int}(\sigma_\infty)\vert\widehat\Pi_{0,3}
\]
whence \deleted{$g\in$}
\[
\operatorname{int}(g)\vert\widehat\Pi_{0,3}=\mathrm{id}
\quad\text{i.e.}\quad
g\in\unsure{\operatorname{Centre}}(\widehat\Pi_{0,3})
\quad\text{i.e.}\quad g=1,
\]
OK.

But under what condition does an automorphism of $\widehat\Pi_{0,3}$ \deleted{can} extend to an automorphism of $\widehat{\mathfrak S}_{0,3}^{+}$ which \deleted{is}\inserted{\unsure{induces}} the identity on $\mathfrak S_{3}$? We see at once that it \uline{is necessary} for this that the outer automorphism of $\widehat\Pi_{0,3}$ which it induces commute with the action of $\mathfrak S_{3}$.

\scanpage{67}
\mspage{\unsure{322|332}}

This is sufficient, for let $\Gamma\inserted{{}\subset\mathrm{Aut\,ext}(\widehat{\Pi}_{0,3})}$ be the group of outer
automorphisms of $\widehat{\Pi}_{0,3}$ which have this property (we are not
concerned for the moment with the
loop structure of $\widehat{\Pi}_{0,3}$), we have
$\Gamma\cap\mathfrak S_3=(1)$ (identifying $\mathfrak S_3$ with a
\unsure{subgroup} of $\mathrm{Aut\,ext}\,\widehat{\Pi}_{0,3}$) for \deleted{\unsure{$\mathfrak S$}} the center
of $\mathfrak S_3$ is trivial, hence
\margin{this is \unsure{evident|evidently}}
\[
\deleted{\unsure{\Gamma\times}}\ \mathfrak S_3\times\Gamma\hookrightarrow\mathrm{Autext}(\widehat{\Pi}_{0,3}),
\]
and we thus have an extension structure
\[
1\unsure{\longrightarrow}\widehat{\Pi}_{0,3}\longrightarrow\mathcal E\longrightarrow\mathfrak S_3\times\Gamma\longrightarrow1
\tag{33}
\]
(for the center of $\widehat{\Pi}_{0,3}$ is trivial), or again
an extension structure
\[
1\longrightarrow\widehat{\mathfrak S}_{0,3}^{+}\unsure{\longrightarrow}\mathcal E\longrightarrow\Gamma\longrightarrow1
\tag{34}
\]
which defines an outer action of $\Gamma$
on $\widehat{\mathfrak S}_{0,3}^{+}$ --- which is the outer action
sought. Thus

\uline{Proposition} the group of outer automorphisms of $\widehat{\Pi}_{0,3}$ which
\unsure{commute} with $\mathfrak S_3$ is canonically isomorphic to the group
of outer automorphisms of $\widehat{\mathfrak S}_{0,3}^{+}$ which induce
on $\mathfrak S_3$ the trivial outer automorphism.
\margin{But \uline{every} automorphism of $\mathfrak S_3$ is inner!}

\scanpage{68}

\uline{Corollary} the group of automorphisms of $\widehat{\Pi}_{0,3}$ which
commute as outer actions with $\mathfrak S_3$, is isomorphic
to the group of automorphisms of $\widehat{\mathfrak S}_{0,3}^{+}$ which
induce on $\mathfrak S_3$ the trivial automorphism (by
restriction of the latter)

\uuline{Corollary}
\[
\diagram{\begin{tikzcd}[ampersand replacement=\&,column sep=small,row sep=large]
1 \arrow[r,no head]
\& \widehat{\Pi}_{0,3} \arrow[r]
\& \mathrm{Aut}\left(\widehat{\Pi}_{0,3}\,\middle|\,
  \begin{array}{l}\mathrm{comm.}\\\mathrm{ext.}^{\unsure{t}}\\\text{with }\mathfrak S_3\end{array}\right) \arrow[r]
\& \mathrm{Autext}\left(\widehat{\Pi}_{0,3}\,\middle|\,
  \begin{array}{l}\mathrm{comm.}\\\text{with }\mathfrak S_3\end{array}\right) \arrow[r]
\& 1
\\
1 \arrow[r,no head]
\& \widehat{\Pi}_{0,3} \arrow[r,no head] \arrow[u,phantom,"{\text{''}}" description]
\& \underbrace{\mathrm{Aut}\left(\widehat{\mathfrak S}_{0,3}^{+}\,\middle|\,
  \begin{array}{l}\text{identity on the}\\\text{quotient }\mathfrak S_3\end{array}\right)}_{\mathcal E^{!}}
  \arrow[r] \arrow[u,"{\sim\ \text{\unsure{restriction} to }\widehat{\Pi}_{0,3}}"]
\& \mathrm{Autext}\left(\widehat{\mathfrak S}_{0,3}^{+}\,\middle|\,
  \begin{array}{l}\text{\deleted{\unsure{finite}}\inserted{\unsure{trivial}}}\\\text{on the}\\\text{quotient }\mathfrak S_3\end{array}\right)
  \arrow[r] \arrow[u,"\sim"]
\& \unsure{1}
\end{tikzcd}}
\tag{35}
\]

\deleted{\uline{Corollary}
Let $u\in\mathrm{Aut}\deleted{\mathrm{ext}}(\widehat{\mathfrak S}_{0,3}^{+})$, \unsure{inducing}
\unsure{Let} \illegible{1} $u\in\mathrm{Aut}(\widehat{\mathfrak S}_{0,3}^{+})$ and $\unsure{u_0}\unsure{=|:=}u|\widehat{\Pi}_{0,3}$
\margin{in $\mathcal E$}
the trivial \inserted{outer} automorphism on $\mathfrak S_3$. We thus have
\unsure{commuting} \illegible{2}
\[
u(\sigma_\infty)=g\sigma_\infty\quad\text{with}\quad g\in\widehat{\Pi}_{0,3}\quad\text{\unsure{well} determined},
\]
we have \illegible{1} that $\sigma_\infty$
and $(g\sigma_\infty)^2=1$ i.e. $g\sigma_\infty g\sigma_\infty=1$ i.e. $g\sigma_\infty g$.}

This extension $\mathcal E^{!}$ of $\Gamma$ by $\widehat{\Pi}_{0,3}$ is
split\inserted{\unsure{noncanonically}}, the \unsure{splitting|splitting} subgroup consists
of the automorphisms of $\widehat{\Pi}_{0,3}$ which commute with $\sigma_\infty$,
or again the automorphisms of $\widehat{\mathfrak S}_{0,3}^{+}$ which
\deleted{\unsure{commute}}\inserted{fix} \unsure{with} $\sigma_\infty$ (and \inserted{which fix} $\rho\ \mathrm{mod}\ \widehat{\Pi}_{0,3}$)
\margin{comes from the fact that there \deleted{\illegible{2}} \inserted{\unsure{exists}} \deleted{\illegible{1}} \unsure{a} \illegible{2} \inserted{$\widehat{\Pi}_{0,3}$-\unsure{conjugation}} \illegible{1} of order $2$ above $\sigma_\infty$ \ldots}
\[
\left\{
\begin{aligned}
u(\sigma_\infty)&=\sigma_\infty\\
u(\rho)&=g.\rho
\end{aligned}
\right.
\qquad(g\in\widehat{\Pi}_{0,3})
\tag{36}
\]

\scanpage{69}
\mspage{323}

\unsure{writing} that $u$ is an \unsure{endomorphism} (for we \unsure{recall}
that $\widehat{\mathfrak S}_{0,3}^{+}$ is defined by the generators $\sigma_\infty,\rho$
with the relations
\begin{equation}
\sigma_\infty^2=\rho^3=1
\tag{37}
\end{equation}
\unsure{we find} the condition
\begin{equation}
(g\rho)^3=1\qquad\text{i.e.}\qquad g\,\rho(g)\,\rho^2(g)=1.
\tag{38}
\end{equation}

\margin{moreover it \unsure{is necessary} in order that $\sigma_\infty$ and $g\rho$
\unsure{generate} $\unsure{\widehat\Pi_{0,3}}$, with only relations
$\sigma_\infty^{\unsure{2}}=(g\rho)^3=1$.}

\uline{Corollary}. Let $u\in\mathcal E^{!}$ be an automorphism of
$\widehat{\mathfrak S}_{0,3}^{+}$ which induces the identity on $\mathfrak S_3$.
Equivalent conditions:

a) $\exists\ p\in\widehat{\mathbb Z}^{*}$ such that
$u(\underset{\varepsilon_0}{\uline{\ell'_0}})$ conjugate in
$\widehat{\mathfrak S}_{0,3}^{+}$ to $\ell_0'{}^p$
(this $p$ is then unique\inserted{, and we can conjugate by an \unsure{element} of $\widehat\Pi_{0,3}$}).

b) $\exists\ p\in\widehat{\mathbb Z}^{*}$ such that $u(\ell_0)$ is conjugate in
$\widehat\Pi_{0,3}$ to $\ell_0^p$ (this $p$ is then unique).

c) $u_0\overset{\mathrm{dfn}}{=}u\mid\widehat\Pi_{0,3}$ respects the loop structure of $\widehat\Pi_{0,3}$.

Moreover, the $p$ in a) and b) is the same.

\uline{\unsure{Proof}}. If $u(\ell'_0)=\operatorname{int}(g_0)(\ell_0'{}^p)$, with
$g_0\in\widehat{\mathfrak S}_{0,3}^{+}$,
squaring we \unsure{will have} $u(\ell_0)=\operatorname{int}(g_0)(\ell_0^p)$,
this shows that the image $\dot g_0$ of $g_0$ in $\mathfrak S_3$
\unsure{leaves invariant} $0$ i.e. is \deleted{\illegible{1}} $1$ or $\dot\sigma_0$,
\deleted{therefore} if it is $1$ we do have b), otherwise we \unsure{put|\deleted{know q}}
$g_0=h\sigma_0$ \inserted{($h\in\widehat\Pi_{0,3}$)}
\deleted{\unsure{whence}}
\[
u(\ell_0)=\operatorname{int}(h)\,
\underset{(\sigma_0(\ell_0))^p}{\uline{\sigma_0(\ell_0^p)}}
\]
now we know that $\sigma_0(\ell_0)=\operatorname{int}(a)\ell_0$
\inserted{(with $a\in\widehat\Pi_{0,3}$)},
\inserted{whence $\sigma_0(\ell_0)^p=\operatorname{int}(a)\ell_0^p$},
thus $u(\ell_0)=\operatorname{int}(ha)(\ell_0^p)$, whence b) with the same $p$.

\scanpage{70}

Since we know the uniqueness of $p$ in b), in a) too. If we have
\begin{equation}
u_0(\ell_0)=\operatorname{int}(g_0)(\ell_0^p)
\qquad (g_0\in\widehat\Pi_{0,3})
\tag{39}
\end{equation}
we deduce from this, writing
\[
\deleted{\rho u_0\rho^{-1}(\ell_0)\unsure{=}}
\qquad
\rho u_0\rho^{-1}=\operatorname{int}(g^{-1})u_0
\qquad (\unsure{a|g}\in\widehat\Pi_{0,3})
\]
\[
\deleted{u_0 \ell_1}
\]
\[
u_0(\ell_1)=\operatorname{int}(\unsure{a|g})
\underbrace{\rho\underbrace{u_0\underbrace{\rho^{-1}(\ell_1)}_{\ell_0}}_{
\operatorname{int}(g_0)(\ell_0^p)}}_{
\operatorname{int}(\rho(g_0))(\ell_1^p)}
=\operatorname{int}(g\rho(g_0))(\ell_1^p)
\]
i.e.
\begin{equation}
u_0(\ell_1)=\operatorname{int}(g_1)(\ell_1^p)
\qquad\text{with}\quad g_1=g\rho(g_0)
\tag{40}
\end{equation}
and likewise
\begin{equation}
\unsure{u_0|u}(\ell_\infty)=\operatorname{int}(g_\infty)(\ell_\infty^p)
\qquad\text{with}\quad
\begin{aligned}
g_\infty&=g\rho(g_1)\\
&=g\rho(g)\rho^2(g_0)
\end{aligned}
\tag{41}
\end{equation}
thus b) $\Rightarrow$ c). Finally \unsure{let us prove} c) $\Rightarrow$ a).

We have by hypothesis
\[
\begin{array}{rcl}
u_0(\ell_0)&=&\operatorname{int}(g)\ell_0^p\\
\text{''}&&\text{''}\\
u(\varepsilon_0)^2&&
\operatorname{int}(g)(\unsure{\ell|\varepsilon}_0^{\unsure{2p|p}})^2
=(\operatorname{int}g(\unsure{\ell_0^p|\varepsilon_0^p}))^2
\end{array}
\]
I say that we then have
\[
u(\varepsilon_0)=\operatorname{int}(g)(\varepsilon_0^p).
\]
It suffices to \deleted{see}\inserted{\unsure{recall}} that
\deleted{\unsure{every}}\inserted{the only} \unsure{element} of
$\widehat{\mathfrak S}_{0,3}^{+}$ with square equal to $\ell_0$ is $\varepsilon_0$\ldots\quad OK

\margin{\unsure{check}}